\section{穿透结果：54只元数据、28只优先池、56份原始PDF}

本轮对54只公司逐一检查券商研报元数据，并对28只核心、核心候选和预告候选各归档最近两份原始PDF及文本，共56份、下载失败0份。以下不是标题目录，而是逐票比较两份报告的预测分母、核心主线、风险与预告时序。若报告早于最新预告，原预测只作为“旧市场预期”，不得继续作为当前估值分母。

\begin{exhibitbox}[Exhibit 21：研报时序审计]
\begin{tabularx}{\textwidth}{L{3.2cm}R{2.0cm}R{2.0cm}X}
\toprule
\textbf{类型} & \textbf{标的数} & \textbf{PDF数} & \textbf{估值后果} \\
\midrule
完整优先池 & 28 & 56 & 每只两份原始PDF，保留预测、风险与本地路径 \\
预告后新研报 & 3 & 4+2页 & 浪潮信息、锐捷网络、工业富联可使用新分母；工业富联另有两份完整报告页 \\
预告前旧研报 & 其余预告公司 & 已归档 & 只作旧预期对照，必须进入重估或等待池 \\
无公开点目标 & 若干 & --- & 不以评级或内部计算冒充Street目标，目标权重为0 \\
\bottomrule
\end{tabularx}
\end{exhibitbox}

\section{香农芯创（300475）：两份原始研报对照}

\textbf{华鑫证券，2026-03-22，《公司动态研究报告：企业级存储需求旺盛，“海普存储”实现年度盈利》}
元数据给出2026E EPS 2.36元、PE 66.5倍；时序为\textbf{预告前，分母可能失效}。原文盈利段摘要：双剑合璧》2024-06-05 预测公司 2025-2027 年收入分别为 352.73、462.89、570.14 公 2、《香农芯创（300475）：2024Q1 亿元，EPS 分别为 1.17、2.36、3.47 元， 求，自主品牌“海普存储”实现年度盈利， 研 上调至“买入”投资评级。 究 风险提示 宏观经济的风险，产品研发不及预期的风险，行业竞争加剧 证券研究报告 的风险，下游需求不及预期的风险。 预测指标 2024A 2025E 2026E 2……风险段摘要：宏观经济的风险，产品研发不及预期的风险，行业竞争加剧 证券研究报告 的风险，下游需求不及预期的风险。 预测指标 2024A 2025E 2026E 2027E 主营收入（百万元） 24,271 35,273 46,289 57,014 增长率（\%） 115.4\% 45.3\% 31.2\% 2……

\textbf{中邮证券，2026-02-05，《海普存储开启年度规模盈利》}
元数据给出2026E EPS 2.36元、PE 64.9倍；时序为\textbf{预告前，分母可能失效}。原文盈利段摘要：和财务指标。风险段摘要：供应商依赖风险，重要产品线的授权取消或不能续约的风险，客 户集中度高的风险，存货保管风险，存货跌价风险，存储器等 IC 产 品价格波动导致毛利率波动的风险，商誉减值的风险。 n 盈利预测和财务指标。

\begin{keyinsight}[AStock处理]
处置为\textbf{预告后重估池}。两份报告均早于最新预告，只能作为旧市场预期；旧EPS不得继续直接入模。下一步门槛：重建全年毛利、海普存储收入与现金流；旧全年EPS 2.36元不得继续使用。
\end{keyinsight}

\section{江波龙（301308）：两份原始研报对照}

\textbf{国信证券，2026-05-07，《1Q26归母净利润同比增长2644.05\%，端侧应用多维拓展》}
元数据给出2026E EPS 26.47元、PE 15.2倍；时序为\textbf{预告前，分母可能失效}。原文盈利段摘要：和财务指标 2024 2025 2026E 2027E 2028E 营业收入(百万元) 17,464 22,766 40,455 48,461 55,012 (+ / -\%) 72.5\% 30.4\% 77.7\% 19.8\% 13.5\% 净利润(百万元) 499 1423 11097 11311 11458 (+ / -\%) 160.2\% 185.4\% 679.6\% 1.9\% 1.3\% 每股收益（元） 1.20 3.40 26.47 26.98 27.34……风险段摘要：存储行业需求不及预期，新产品放量不及预期等。 盈利预测和财务指标 2024 2025 2026E 2027E 2028E 营业收入(百万元) 17,464 22,766 40,455 48,461 55,012 (+ / -\%) 72.5\% 30.4\% 77.7\% 19.8\% 13.5\%……

\textbf{爱建证券，2026-05-07，《2025年报\&2026Q1点评：国产存储模组龙头进入业绩爆发期》}
元数据给出2026E EPS 47.64元、PE 10.1倍；时序为\textbf{预告前，分母可能失效}。原文盈利段摘要：报告期 2024 2025 2026E 2027E 2028E 20.6\% 朱俊宇 归母净利润（百万元） 499 1,423 19,969 21,908 22,740 同比增长（\%） -160.2\% 185.4\% 1303.0\% 9.7\% 3.8\% 021-32229888-25520 每股收益（元 / 股） 1.19 3.40 47.64 52.27 54.25 毛利率（\%） 19.0\% 19.4\% 64.6\% 57.2\% 49.8\% ROE（\%） 7.……风险段摘要：1)客户集中风险 2)新业务商业化进度不及预期风险 3)下游需求不及预期。 0755-83562506 财务数据及盈利预测 报告期 2024 2025 2026E 2027E 2028E 20.6\% 朱俊宇 归母净利润（百万元） 499 1,423 19,969 21,908 22,740……

\begin{keyinsight}[AStock处理]
处置为\textbf{正式核心模型}。两份报告均早于最新预告，只能作为旧市场预期；旧EPS不得继续直接入模。下一步门槛：H2归母净利不低于80亿元、经营现金流改善、无重大存货减值。
\end{keyinsight}

\section{兆易创新（603986）：两份原始研报对照}

\textbf{华鑫证券，2026-06-29，《公司动态研究报告：存储、MCU迎涨价浪潮，多业务推动业绩强劲增长》}
元数据给出2026E EPS 8.35元、PE 92.2倍；时序为\textbf{预告前，分母可能失效}。原文盈利段摘要：研 预测公司 2026-2028 年收入分别为 170.48、213.70、262.88 究 亿元，EPS 分别为 8.35、10.55、12.68 元， 买入”投资评级。 风险提示 宏观经济的风险，产品研发不及预期的风险，行业竞争加剧 的风险，下游需求不及预期的风险。 预测指标 2025A 2026E 2027E 2028E 主营收入（百万元） 9,203 17,048 21,370 26,288 增长率（\%） 25.1\% 85.2\% 25.4\% 23.……风险段摘要：宏观经济的风险，产品研发不及预期的风险，行业竞争加剧 的风险，下游需求不及预期的风险。 预测指标 2025A 2026E 2027E 2028E 主营收入（百万元） 9,203 17,048 21,370 26,288 增长率（\%） 25.1\% 85.2\% 25.4\% 23.0\% 归母净利……

\textbf{中邮证券，2026-05-27，《全芯赋能，智创未来》}
元数据给出2026E EPS 8.67元、PE 59.5倍；时序为\textbf{预告前，分母可能失效}。原文盈利段摘要：和财务指标 项目 / 年度 2025A 2026E 2027E 2028E 营业收入（百万元） 9203 16923 21112 25403 增长率(\%) 25.12 83.87 24.75 20.32 EBITDA（百万元） 1950.75 7331.47 8809.84 10346.05 归属母公司净利润（百万元） 1648.02 6081.03 7458.85 9045.60 增长率(\%) 49.47 268.99 22.66 21.27 EPS(元……风险段摘要：宏观环境和行业波动风险，供应链风险，人才流失风险，汇兑损 益风险，产品研发不及预期的风险，外延并购整合不及预期的风险。 n 盈利预测和财务指标 项目 / 年度 2025A 2026E 2027E 2028E 营业收入（百万元） 9203 16923 21112 25403 增长率(\%) 2……

\begin{keyinsight}[AStock处理]
处置为\textbf{预告后重估池}。两份报告均早于最新预告，只能作为旧市场预期；旧EPS不得继续直接入模。下一步门槛：拆分证券公允价值收益与扣非利润，重建存储和MCU全年量价假设。
\end{keyinsight}

\section{中兴通讯（000063）：两份原始研报对照}

\textbf{国金证券，2026-05-05，《连接+算力双轮驱动，利润筑底回升》}
元数据给出2026E EPS 1.54元、PE 24.3倍；时序为\textbf{无预告时序}。原文盈利段摘要：归母净利润增长率 -9.66\% -33.32\% 31.24\% 25.23\% 6.97\% 摊薄每股收益(元) 1.761 1.174 1.541 1.930 2.065 我 们 预 计 公 司 2026-2028 年 营 收 分 别 为 每股经营性现金流净额 2.59 0.57 1.80 3.23 3.03 1728.08 / 2129.83 / 2500.84 亿元，同比增长 29.1\% / 23.3\% / 17.4\%； ROE(归属母公司)(摊薄) 11……风险段摘要：5G 发展不及预期风险；中美贸易摩擦加剧风险；AI 业务发展不及 预期。 敬请参阅最后一页特别声明 1 扫码获取更多服务 公司点评 附录：三张报表预测摘要 损益表（人民币百万元） 资产负债表（人民币百万元） 2023 2024 2025 2026E 2027E 2028E 2023 202……

\textbf{开源证券，2026-03-10，《公司信息更新报告：算力业务跨越式增长，研发投入夯实长期竞争力》}
元数据给出2026E EPS 1.44元、PE 21.0倍；时序为\textbf{无预告时序}。原文盈利段摘要：预计 2026-2028 年公司归母净利润分别为 68.9、82.5、95.9 亿元，EPS 分别为 1.44、1.73、2.01 元， 持续建设，我们看好第二曲线业务长期发展，维持“买 中兴通讯 沪深300 入”评级。 60\% 40\% 算力业务跨越式增长，消费者业务聚焦全场景 AI 终端布局 20\% 分业务板块看，公司运营商网络业务实现营收 628.6 亿元，占总营收 46.9\%，其 0\% 中，国内市场受通信基础设施投资下降影响，营收承压；公司政企业务全……风险段摘要：网络建设不及预期、政企客户拓展不及预期、中美贸易摩擦加剧。 《发布智算超节点服务器，第二曲线 财务摘要和估值指标 蓄势增长—公司信息更新报告》 指标 2024A 2025A 2026E 2027E 2028E -2025.4.29 营业收入(百万元) 121,299 133,895 152……

\begin{keyinsight}[AStock处理]
处置为\textbf{正式核心模型}。公司未形成H1预告时序，本轮报告用于Q1和全年预测对照，仍需按当前价格重算。下一步门槛：算力收入占比不低于25\%、H2利润恢复正增长、经营现金流改善。
\end{keyinsight}

\section{紫光股份（000938）：两份原始研报对照}

\textbf{华鑫证券，2026-05-05，《公司事件点评报告：公司业绩高增，AI解决方案领导者深度受益于算力红利》}
元数据给出2026E EPS 0.85元、PE 37.6倍；时序为\textbf{预告前，分母可能失效}。原文盈利段摘要：预 测 公 司 2026-2028 年 收 入 分 别 为 1147.15 、 1332.29 、 1510.95 亿元，EPS 分别为 0.85、1.10、1.43 元， 风险提示 宏观经济波动导致政企及运营商 IT 资本开支不及预期； AI 相关技术迭代及新产品研发进展不及预期；核心零部件供应 链面临外部环境限制 等风险。 预测指标 2025A 2026E 2027E 2028E 主营收入（百万元） 96,748 114,715 133,229 151……风险段摘要：宏观经济波动导致政企及运营商 IT 资本开支不及预期； AI 相关技术迭代及新产品研发进展不及预期；核心零部件供应 链面临外部环境限制 等风险。 预测指标 2025A 2026E 2027E 2028E 主营收入（百万元） 96,748 114,715 133,229 151,095 增长……

\textbf{国金证券，2026-05-05，《新华三增收增利，积极备货支撑全年交付》}
元数据给出2026E EPS 0.87元、PE 31.5倍；时序为\textbf{预告前，分母可能失效}。原文盈利段摘要：每股经营性现金流净额 0.85 1.22 -2.87 2.74 0.01 我们预计公司 2026-2028 年营收分别为 1145 / 1325 / 1523 亿元， ROE(归属母公司)(摊薄) 11.79\% 11.43\% 14.67\% 15.73\% 17.02\% 同比增长 18\% / 16\% / 15\%；预计净利润分别为 25.0 / 31.2 / 40.1 亿 P / E 50.62 41.74 31.50 25.18 19.59 P / B 5.9……风险段摘要：存货减值风险、芯片供应短缺风险、市场竞争加剧、Scale-up 域 研发进展不及预期。 敬请参阅最后一页特别声明 1 扫码获取更多服务 公司点评 附录：三张报表预测摘要 损益表（人民币百万元） 资产负债表（人民币百万元） 2023 2024 2025 2026E 2027E 2028E 2……

\begin{keyinsight}[AStock处理]
处置为\textbf{预告后重估池}。两份报告均早于最新预告，只能作为旧市场预期；旧EPS不得继续直接入模。下一步门槛：扣除2.5--3.5亿元一次性收益，重建新华三持股比例与利息费用后的全年EPS。
\end{keyinsight}

\section{浪潮信息（000977）：两份原始研报对照}

\textbf{中邮证券，2026-07-09，《深度布局国产算力生态，业绩表现大超市场预期》}
元数据给出2026E EPS 2.53元、PE 28.1倍；时序为\textbf{预告后}。原文盈利段摘要：和财务指标 项目 / 年度 2025A 2026E 2027E 2028E 营业收入（百万元） 164782 212141 253413 290773 增长率(\%) 43.25 28.74 19.46 14.74 EBITDA（百万元） 3321.94 4552.24 5332.08 6046.38 归属母公司净利润（百万元） 2412.87 3719.02 4467.68 5192.31 增长率(\%) 5.20 54.13 20.13 16.22 EPS……风险段摘要：市场竞争加剧风险；技术更新换代过快风险；下游需求波动风险； 供应链安全风险等。

\textbf{群益证券，2026-07-08，《二季度业绩超预期增长，AI服务器有望持续抬高盈利能力》}
元数据给出2026E EPS 3.86元、PE 18.0倍；时序为\textbf{预告同日}。原文盈利段摘要：我们预计公司 2026-2028 年净利润分别为 56.67 / 60.61 / 72.65 亿 元，YOY 分别为+134.9\% / +7.0\% / +19.9\%；EPS 分别为 3.86 / 4.13 / 4.95 元，当 前股价对应 A 股 2026-2028 年 P / E 为 18 / 17 / 14 倍，维持“买进”建议 风险提示：1、传统服务器行业景气度下滑；2、美国相关芯片供应政策 变化；3、下游云计算、AI 产业发展不及预期。 年度截止……风险段摘要：1、传统服务器行业景气度下滑；2、美国相关芯片供应政策 变化；3、下游云计算、AI 产业发展不及预期。 年度截止 12 月 31 日 2024 2025 2026E 2027E 2028E 纯利 (Net profit) RMB 百万元 2291.76 2412.87 5666.84 60……

\begin{keyinsight}[AStock处理]
处置为\textbf{业绩兑现、等待回撤}。报告预测晚于或同步于预告，可用于更新分母，但仍需按7月10日现价重算空间。下一步门槛：预告后EPS已刷新，但现价接近群益90元目标；等待毛利率和营运资金确认。
\end{keyinsight}

\section{智微智能（001339）：两份原始研报对照}

\textbf{中邮证券，2026-06-18，《智算业务兑现增长，LPU与具身智能打开AI硬件新空间》}
元数据给出2026E EPS 2.48元、PE 39.2倍；时序为\textbf{预告前，分母可能失效}。原文盈利段摘要：我们认为，公司智算业务是当前业绩兑现的主线，LPU 与具身智 能是中长期成长期权。预计公司 2026-2028 年实现营收 63.78、78.11、 95.79 亿元，同比增长 56.07\%、22.47\%、22.64\%；实现归母净利润分。风险段摘要：算力租赁价格回落及资产利用率不及预期；智算中心建设和定增 进度不及预期；元川微 LPU 产品研发、验证与商业化不及预期；具身 智能产业量产节奏不及预期；原材料价格和供应链波动风险。 盈利预测和财务指标 项目 / 年度 2025A 2026E 2027E 2028E 营业收入（百万元） 40……

\textbf{太平洋，2026-05-13，《智算业务高速增长，战略布局LPU》}
元数据给出2026E EPS 1.65元、PE 62.9倍；时序为\textbf{预告前，分母可能失效}。原文盈利段摘要：和财务指标 2025A 2026E 2027E 2028E 营业收入（百万元） 4,087 4,801 5,515 6,275 营业收入增长率(\%) 1.30\% 17.48\% 14.87\% 13.77\% 归母净利（百万元） 171 415 597 820 净利润增长率(\%) 36.84\% 142.84\% 43.75\% 37.41\% 摊薄每股收益（元） 0.68 1.65 2.37 3.25 市盈率（PE） 152.72 62.89 43.75 31.84……风险段摘要：下游需求不及预期，市场竞争加剧，供应链波动风险。

\begin{keyinsight}[AStock处理]
处置为\textbf{价格高于旧外部锚}。两份报告均早于最新预告，只能作为旧市场预期；旧EPS不得继续直接入模。下一步门槛：AI算力/ICT出货兑现，同时当前价格回到新盈利预测可解释区间。
\end{keyinsight}

\section{风华高科（000636）：两份原始研报对照}

\textbf{中邮证券，2026-03-08，《持续高景气》}
元数据给出2026E EPS 0.42元、PE 56.1倍；时序为\textbf{预告前，分母可能失效}。原文盈利段摘要：和财务指标 项目 / 年度 2024A 2025E 2026E 2027E 营业收入（百万元） 4939 5784 6880 7863 增长率(\%) 17.00 17.10 18.95 14.29 EBITDA（百万元） 907.32 1136.41 1417.43 1718.10 归属母公司净利润（百万元） 337.37 357.08 489.66 640.33 增长率(\%) 94.47 5.84 37.13 30.77 EPS(元 / 股) 0.29……风险段摘要：行业及市场变动风险；市场竞争加剧风险；供应链成本上行风险； 涨价不及预期风险；产能释放延迟风险；下游需求不及预期风险。

\textbf{中邮证券，2025-11-19，《风华铸器，高科赋能》}
元数据给出2026E EPS 0.45元、PE 38.8倍；时序为\textbf{预告前，分母可能失效}。原文盈利段摘要：和财务指标 项目 / 年度 2024A 2025E 2026E 2027E 营业收入（百万元） 4939 5856 6694 7537 增长率(\%) 17.00 18.56 14.31 12.60 EBITDA（百万元） 907.32 1099.58 1399.76 1712.67 归属母公司净利润（百万元） 337.37 364.79 517.47 676.66 增长率(\%) 94.47 8.13 41.85 30.76 EPS(元 / 股) 0.29……风险段摘要：行业及市场变动风险；市场竞争加剧风险；供应链成本上行风险； 下游需求不及预期风险。

\begin{keyinsight}[AStock处理]
处置为\textbf{周期兑现但价格要求高}。两份报告均早于最新预告，只能作为旧市场预期；旧EPS不得继续直接入模。下一步门槛：Q2利润加速同时MLCC价格、产品结构和现金流持续改善。
\end{keyinsight}

\section{亿道信息（001314）：两份原始研报对照}

\textbf{华鑫证券，2026-07-10，《公司动态研究报告：端侧AI前景广阔，先进封装打造第二增长曲线》}
元数据给出2026E EPS 1.66元、PE 35.0倍；时序为\textbf{预告前，分母可能失效}。原文盈利段摘要：预测公司 2026-2028 年收入分别为 58、74、92 亿元，EPS 分 别为 1.66、1.80、2.33 元， 公司参与投资发起的先进封装项目与端侧 AI 主业形成协同， 有望打造公司在半导体领域的第二增长曲线，同时赋能主业 发展。我们看好公司未来的业务拓展和业绩增长，维持“买 入”投资评级。 风险提示 1）上游产能受限供应不足风险；市场竞争加剧风险；2）下 游需求不及预期风险。 预测指标 2025A 2026E 2027E 2028E 主营收入……风险段摘要：1）上游产能受限供应不足风险；市场竞争加剧风险；2）下 游需求不及预期风险。 预测指标 2025A 2026E 2027E 2028E 主营收入（百万元） 4,583 5,805 7,389 9,186 增长率（\%） 44.2\% 26.6\% 27.3\% 24.3\% 归母净利润（百万元） 6……

\textbf{中邮证券，2025-02-05，《加速创新，智塑未来》}
元数据给出2026E EPS 0.86元、PE 56.6倍；时序为\textbf{预告前，分母可能失效}。原文盈利段摘要：和财务指标 项目 / 年度 2023A 2024E 2025E 2026E 营业收入（百万元） 2594 3203 3804 4503 增长率(\%) -5.82 23.50 18.77 18.35 EBITDA（百万元） 117.86 74.43 111.89 149.85 归属母公司净利润（百万元） 128.81 41.60 80.50 121.61 增长率(\%) -34.67 -67.70 93.50 51.07 EPS(元 / 股) 0.91 0.2……风险段摘要：技术迭代和研发投入不及预期的风险；行业竞争格局加剧风险。

\begin{keyinsight}[AStock处理]
处置为\textbf{扣除一次性后观察}。两份报告均早于最新预告，只能作为旧市场预期；旧EPS不得继续直接入模。下一步门槛：剔除约0.86亿元投资收益，单独验证AI终端主营利润。
\end{keyinsight}

\section{多氟多（002407）：两份原始研报对照}

\textbf{民生证券，2025-04-24，《2024年年报及2025年一季报点评：业绩持续承压，六氟磷酸锂出货稳步增长》}
元数据给出2026E EPS 0.36元、PE 32.0倍；时序为\textbf{预告前，分母可能失效}。原文盈利段摘要：与财务指标 项目 / 年度 2024A 2025E 2026E 2027E 营业收入（百万元） 8,207 10,742 13,408 16,190 增长率（\%） -31.3 30.9 24.8 20.8 归属母公司股东净利润（百万元） -308 290 425 625 增长率（\%） -160.4 194.0 46.5 47.2 每股收益（元） -0.26 0.24 0.36 0.53 PE / 47 32 22 PB 1.6 1.6 1.6 1.5 资料……风险段摘要：宏观经济与市场波动风险，公司规模扩大导致的管理风险，产品 023 / 09 / 04 和技术更新风险，环保及安全生产风险。 盈利预测与财务指标 项目 / 年度 2024A 2025E 2026E 2027E 营业收入（百万元） 8,207 10,742 13,408 16,190 增长率（……

\textbf{华安证券，2024-09-02，《1H24业绩承压下滑，多元化布局协同发展》}
元数据给出2026E EPS 0.60元、PE 16.7倍；时序为\textbf{预告前，分母可能失效}。原文盈利段摘要：资产负债表 单位:百万元 利润表 单位:百万元 会计年度 2023A 2024E 2025E 2026E 会计年度 2023A 2024E 2025E 2026E 流动资产 12169 10107 11029 13713 营业收入 11937 9853 13454 19286 现金 6875 5595 4992 5222 营业成本 9996 8710 11704 16637 应收账款 2081 1374 2111 2858 营业税金及附加 78 61 85……风险段摘要：（1）原材料价格上涨带来成本的升高； （2）客户验证进度不及预期； （3）产能释放进度不及预期； （4）下游需求不及预期。 重要财务指标 单位:百万元 主要财务指标 2023A 2024E 2025E 2026E 营业收入 11937 9853 13454 19286 收入同比（\%） -3……

\begin{keyinsight}[AStock处理]
处置为\textbf{周期观察、Q2减速}。两份报告均早于最新预告，只能作为旧市场预期；旧EPS不得继续直接入模。下一步门槛：六氟磷酸锂价格与销量继续改善，且Q2减速不延续、现金流转正。
\end{keyinsight}

\section{翔鹭钨业（002842）：两份原始研报对照}

\textbf{中邮证券，2018-12-04，《专注钨制品行业，短期业绩弹性明显》}
元数据给出2026E EPS ---元、PE ---倍；时序为\textbf{预告前，分母可能失效}。原文盈利段摘要：与评级：预计公司 2018-2020 年 EPS 分别为 0.62 元、0.58 元和 0.52 元，目前股价对应 2018-2020 年市盈率分别为 27.3 倍、29.3 倍和 32.6 倍，给予“中性”的投资评级。 风险提示：经济下行，下游需求大幅下滑；产能投放进度晚于预期； 汇率波动风险。 会计年度 2017A 2018E 2019E 2020E 营业收入（百万元） 976 1,621 1,701 1,616 增长率（\%） 36.5\% 66.1\%……风险段摘要：经济下行，下游需求大幅下滑；产能投放进度晚于预期； 汇率波动风险。 会计年度 2017A 2018E 2019E 2020E 营业收入（百万元） 976 1,621 1,701 1,616 增长率（\%） 36.5\% 66.1\% 5.0\% -5.0\% 净利润（百万元） 69 105 98 8……

\textbf{太平洋，2018-10-24，《2018年三季报点评：业绩再创新高，产能扩张或对冲价格下行压力》}
元数据给出2026E EPS ---元、PE ---倍；时序为\textbf{预告前，分母可能失效}。原文盈利段摘要：和财务指标： 2017A 2018E 2019E 2020E 执业资格证书编码： 营业收入(百万元) 976 1763 2011 2295 （+ / -\%） 36.5 80.6 14.1 14.1 净利润(百万元） 69 126 149 181 （+ / -\%） 20.61\% 83.49\% 17.77\% 21.21\% 摊薄每股收益(元) 0.70 1.26 1.49 1.80 市盈率(PE) 62.61 14.94 12.68 10.46 资料来源：Win……风险段摘要：贸易战、需求不及预期、泛亚库存释放。 E-MAIL： 盈利预测和财务指标： 2017A 2018E 2019E 2020E 执业资格证书编码： 营业收入(百万元) 976 1763 2011 2295 （+ / -\%） 36.5 80.6 14.1 14.1 净利润(百万元） 69 126……

\begin{keyinsight}[AStock处理]
处置为\textbf{周期观察、现金流门槛}。两份报告均早于最新预告，只能作为旧市场预期；旧EPS不得继续直接入模。下一步门槛：钨价可顺利传导，下游需求不被高价压制，经营现金流覆盖利润。
\end{keyinsight}

\section{方正科技（600601）：两份原始研报对照}

\textbf{中航证券，2025-07-24，《锚定AI算力浪潮，高端PCB引领增长新势能》}
元数据给出2026E EPS 0.11元、PE 50.1倍；时序为\textbf{预告前，分母可能失效}。原文盈利段摘要：原文未形成可稳定抽取的连续段落，保留标题、预测表与PDF路径供复核。风险段摘要：原文未形成可稳定抽取的连续段落，保留标题、预测表与PDF路径供复核。

\textbf{上海证券，2020-02-25，《方正集团债务重组 公司估值加速修复》}
元数据给出2026E EPS ---元、PE ---倍；时序为\textbf{预告前，分母可能失效}。原文盈利段摘要：我们预期公司 2019-2021 年实现营业收入 59.86 亿元、65.53 亿元、68.51 亿元，同比增长分别为 5.00\%、6.13\%和 7.84\%；归属于母公司股东净 利润为-13.01 亿元、0.26 亿元和 2.51 亿元；EPS 分别为-0.59 元、0.01 元和 0.11 元，对应 2021 年 PE 为 30X。未来六个月内，首次覆盖给与 “增持”评级。 主要风险 （1）股权架构梳理不及预期；（2）宽带业务降低盈利能力。 报告编号：……风险段摘要：6 七、 附表 7 图 图 1 方正科技是方正集团旗下上市公司 …。

\begin{keyinsight}[AStock处理]
处置为\textbf{预告后重估池}。两份报告均早于最新预告，只能作为旧市场预期；旧EPS不得继续直接入模。下一步门槛：重建高端PCB收入、毛利、客户订单和资本开支，旧EPS 0.11元作废。
\end{keyinsight}

\section{恺英网络（002517）：两份原始研报对照}

\textbf{中邮证券，2026-06-08，《业绩稳健增长，传奇盒子打造第二增长曲线》}
元数据给出2026E EPS 1.24元、PE 12.5倍；时序为\textbf{无预告时序}。原文盈利段摘要：我们看好公司传奇盒子快速放量，构筑第二增长曲线，同时以前 瞻性 AI 布局卡位游戏行业未来方向。我们预测公司 2026-2028 年实 现营收 78.22、85.42、95.88 亿元，同比+46.9\%、9.2\%、12.24\%；实 现归母净利润 26.46、30.48、36.28 亿元，同比+39.01\%、15.20\%、 19.02\%。维持“买入”评级。 风险提示： 行业监管风险、新业务不及预期、新游推广不及预期等。 盈利预测和财务指标 项目 / 年度 2……风险段摘要：行业监管风险、新业务不及预期、新游推广不及预期等。 盈利预测和财务指标 项目 / 年度 2025A 2026E 2027E 2028E 营业收入（百万元） 5325 7822 8542 9588 增长率(\%) 4.04 46.90 9.20 12.24 EBITDA（百万元） 2039.8……

\textbf{太平洋，2026-05-05，《游戏盒子驱动营收增长，新游储备及AI布局逐步落地》}
元数据给出2026E EPS 1.20元、PE 14.0倍；时序为\textbf{无预告时序}。原文盈利段摘要：执业资格证书编码： 公司多款老游戏长线运营稳健，贡献流水基本盘；游戏盒子贡献业绩增量。 未来，丰富的游戏储备、以及多元的 AI 布局亦有望贡献增量。因此，我 证券 .22\%，归母净利润分别为 25.7 / 30.4 / 34.8 亿元，对应增速 执业资格证书编码： 35.2\% / 18.1\% / 14.6\%。维持公司“买入”评级。风险段摘要：政策监管趋紧的风险，游戏行业增速放缓的风险，游戏流水、上线节奏不 及预期的风险。 盈利预测和财务指标 2025 2026E 2027E 2028E 营业收入（百万元） 5325 7553 8385 9159 营业收入增长率(\%) 4.04\% 41.83\% 11.02\% 9.22\% 归母净利……

\begin{keyinsight}[AStock处理]
处置为\textbf{券商支持、等待中报}。公司未形成H1预告时序，本轮报告用于Q1和全年预测对照，仍需按当前价格重算。下一步门槛：等待新品流水、用户付费、内容成本和AI收入对EPS的贡献。
\end{keyinsight}

\section{三七互娱（002555）：两份原始研报对照}

\textbf{太平洋，2026-05-05，《业绩稳健增长，多款SLG储备有望开启新品周期》}
元数据给出2026E EPS 1.43元、PE 15.0倍；时序为\textbf{无预告时序}。原文盈利段摘要：公司存量游戏长线运营稳健，且多款新游戏贡献增量。未来，丰富的 游戏储备有望推动新产品周期开启，关注后续新游上线节奏及表现。 因此，我们预计公司 2026-2028 年营收分别为 178 / 192 / 203 亿元，对 应增速 11.6\% / 7.76\% / 5.84\%，归母净利润分别为 31.7 / 34.2 / 36.6 亿。风险段摘要：政策监管趋紧的风险，游戏行业增速放缓的风险，游戏流水、上线节 奏不及预期的风险。 盈利预测和财务指标 2025 2026E 2027E 2028E 营业收入（百万元） 15966 17815 19197 20319 营业收入增长率(\%) -8.46\% 11.58\% 7.76\% 5.84\%……

\textbf{东吴证券，2026-04-22，《2025年报点评：业绩低于预期，关注新游上线进展》}
元数据给出2026E EPS 1.43元、PE 15.3倍；时序为\textbf{无预告时序}。原文盈利段摘要：2024A 2025A 2026E 2027E 2028E 营业总收入（百万元） 17,441 15,966 17,934 19,241 20,561 股价走势 同比（\%） 5.40 (8.46) 12.33 7.29 6.86 三七互娱 沪深300 归母净利润（百万元） 2,673 2,900 3,171 3,405 3,686 92\% 同比（\%） 0.54 8.50 9.33 7.39 8.26 81\% 70\% 59\% EPS-最新摊薄（元 / 股）……风险段摘要：新游表现不及预期，买量竞争加剧，行业监管风险。 1 / 3 东吴证券研究所。

\begin{keyinsight}[AStock处理]
处置为\textbf{券商支持、等待中报}。公司未形成H1预告时序，本轮报告用于Q1和全年预测对照，仍需按当前价格重算。下一步门槛：等待新品流水、用户付费、内容成本和AI收入对EPS的贡献。
\end{keyinsight}

\section{通富微电（002156）：两份原始研报对照}

\textbf{开源证券，2026-05-06，《公司信息更新报告：业绩重回增长快车道，资本开支开启新周期》}
元数据给出2026E EPS 0.98元、PE 52.8倍；时序为\textbf{无预告时序}。原文盈利段摘要：预计 2026-2028 年归母净利润为 14.89 / 18.51 / 23.41 报 亿元（2026-2027 原值 15.95 / 21.31 亿元） ， 化封测业务突 股价走势图 破”三重逻辑下的成长机会，维持“买入”评级。 3nm 多芯片产品封装已通过验证，高端化突破打开业务增长空间 通富微电 沪深300 2025 年通富苏州营收 79.71 亿元 / yoy+3.87\%，通富槟城营收 94.11 亿元 / yoy 160\% +23.08\%，槟城……风险段摘要：研发及量产进度不及预期、半导体周期波动、行业竞争加剧。 财务摘要和估值指标 指标 2024A 2025A 2026E 2027E 2028E 营业收入(百万元) 23,882 27,921 33,070 37,640 43,174 YOY(\%) 7.2 16.9 18.4 13.8 14.……

\textbf{国信证券，2026-04-21，《2025年归母净利润增长80\%，四季度收入创季度新高》}
元数据给出2026E EPS 0.91元、PE 52.9倍；时序为\textbf{无预告时序}。原文盈利段摘要：和财务指标 2024 2025 2026E 2027E 2028E 《通富微电（002156.SZ）-一季度收入同比增长 13.8\%，拟收购 营业收入(百万元) 23,882 27,921 32,329 36,005 39,572 京隆科技 26\%股权》 ——2024-04-28 (+ / -\%) 7.2\% 16.9\% 15.8\% 11.4\% 9.9\% 归母净利润(百万元) 678 1219 1383 1694 1964 (+ / -\%) 299.9\% 7……风险段摘要：新产品研发不及预期；需求不及预期；大客户集中的风险。 二期项目扩产顺利》 ——2024-11-01 《通富微电（002156.SZ）-上半年收入同比增长 11.8\%，盈利能 力提高》 ——2024-08-31 盈利预测和财务指标 2024 2025 2026E 2027E 2028E 《通……

\begin{keyinsight}[AStock处理]
处置为\textbf{价格先行、等待业绩}。公司未形成H1预告时序，本轮报告用于Q1和全年预测对照，仍需按当前价格重算。下一步门槛：等待H1实际利润、先进封装收入、产能利用率和资本开支回报。
\end{keyinsight}

\section{华天科技（002185）：两份原始研报对照}

\textbf{华龙证券，2026-04-29，《公司深度研究：聚焦先进封装，迈向全球领先封测企业》}
元数据给出2026E EPS 0.30元、PE 42.4倍；时序为\textbf{无预告时序}。原文盈利段摘要：及投资评级：半导体景气周期仍在高位，中国大陆封测 产业升级与国产替代不断提速。根据公司公告，2026 年营业收入 目标为 200 亿元。具体来看，公司 2025 年扩产节奏进一步加快， 尤其是在先进封装方面：华天先进进口线已于 2025 年二季度跑 通、国产线于下半年完成搭建和调试并进行小批量样品测试；盘 古半导体于 2025 年部分投产，其他产能将于 2026-2028 年陆续投 产。若干先进封装产能的陆续建成投产，不仅提升公司收入规模， 其高附加值属性……风险段摘要：下游需求疲软；地缘政治与政策风险；国产替代及研 发进度不及预期；收购事项存在不确定性；宏观经济恢复不及预 期；数据引用风险。 盈利预测简表 预测指标 2024A 2025A 2026E 2027E 2028E 营业收入（百万元） 14,462 17,214 20,742 25,824 29……

\textbf{华金证券，2025-08-22，《25H1订单/业绩稳健增长，持续加码2.5D/3D封测》}
元数据给出2026E EPS 0.37元、PE 29.6倍；时序为\textbf{无预告时序}。原文盈利段摘要：我们预计 2025-2027 年，公司营收分别为 175.61 / 206.65 / 232.38 亿元， 同比分别为 21.4\% / 17.7\% / 12.5\%；归母净利润分别为 8.56 / 11.87 / 15.27 亿元，同 比分别为 38.9\% / 38.7\% / 28.7\%；对应 PE 为 41.1 / 29.6 / 23.0 倍。随着人工智能发展 对算力芯片需求加剧，有望带动先进封装需求，考虑到华天科技基于 3D Matrix 平 台，通过集……风险段摘要：行业与市场波动风险；国际贸易摩擦风险；新技术、新工艺、新产品无 法如期产业化风险；扩产不及预期风险；主要原材料 / 设备供应及价格变动风险； http: / / www.huajinsc.cn / 1 / 5。

\begin{keyinsight}[AStock处理]
处置为\textbf{价格先行、等待业绩}。公司未形成H1预告时序，本轮报告用于Q1和全年预测对照，仍需按当前价格重算。下一步门槛：等待H1实际利润、先进封装收入、产能利用率和资本开支回报。
\end{keyinsight}

\section{长电科技（600584）：两份原始研报对照}

\textbf{中邮证券，2026-07-03，《先进封装，芯连万象》}
元数据给出2026E EPS 1.15元、PE 83.6倍；时序为\textbf{无预告时序}。原文盈利段摘要：和财务指标 项目 / 年度 2025A 2026E 2027E 2028E 营业收入（百万元） 38871 43371 49033 54596 增长率(\%) 8.09 11.57 13.06 11.35 EBITDA（百万元） 6132.76 7571.04 9147.44 10800.93 归属母公司净利润（百万元） 1565.24 2053.52 2640.95 3266.59 增长率(\%) -2.75 31.20 28.61 23.69 EPS(元……风险段摘要：行业波动风险；产业政策变化风险；市场竞争加剧风险；汇率波 动风险。 n 盈利预测和财务指标 项目 / 年度 2025A 2026E 2027E 2028E 营业收入（百万元） 38871 43371 49033 54596 增长率(\%) 8.09 11.57 13.06 11.35 EBI……

\textbf{华鑫证券，2026-05-12，《公司事件点评报告：盈利能力复苏，先进封装龙头受益于AI算力强劲需求》}
元数据给出2026E EPS 1.08元、PE 51.8倍；时序为\textbf{无预告时序}。原文盈利段摘要：预测公司 2026-2028 年收入分别为 439.44、492.83、541.2 亿 元，EPS 分别为 1.08、1.37、1.61 元， 投资评级。 风险提示 宏观经济波动；先进封装技术研发与新产能量产导入不及预 期；汇率波动风险。 请阅读最后一页重要。风险段摘要：宏观经济波动；先进封装技术研发与新产能量产导入不及预 期；汇率波动风险。 请阅读最后一页重要。

\begin{keyinsight}[AStock处理]
处置为\textbf{价格先行、等待业绩}。公司未形成H1预告时序，本轮报告用于Q1和全年预测对照，仍需按当前价格重算。下一步门槛：等待H1实际利润、先进封装收入、产能利用率和资本开支回报。
\end{keyinsight}

\section{恒瑞医药（600276）：两份原始研报对照}

\textbf{西南证券，2026-07-01，《研发能力再得验证，创新成果加速落地》}
元数据给出2026E EPS 1.44元、PE 36.8倍；时序为\textbf{无预告时序}。原文盈利段摘要：预计公司 2026-2028 年归母净利润分别为 95 / 112 / 133 亿 基 础数据 元，对应 PE 为 37X / 31X / 26X。公司创新药管线持续布局，国际化出海成果颇 总股本(亿股) 66.37 丰，龙头效应明显，维持“买入”评级。 流通 A 股(亿股) 63.79 52 周内股价区间(元) 45.94-72.47 风险提示：研发及商业化不及预期风险；行业政策风险；国际地缘政治风险等。 总市值(亿元) 3,513.07 总资产(亿元)……风险段摘要：研发及商业化不及预期风险；行业政策风险；国际地缘政治风险等。 总市值(亿元) 3,513.07 总资产(亿元) 712.94 指标 / 年度 2025A 2026E 2027E 2028E 每股净资产(元) 9.59 营业收入（百万元） 31629.42 35770.38 40508.27……

\textbf{华源证券，2026-05-07，《创新成果密集落地，全球化布局加速推进》}
元数据给出2026E EPS 1.45元、PE 37.0倍；时序为\textbf{无预告时序}。原文盈利段摘要：（人民币） 2024 2025 2026E 2027E 2028E 营业收入（百万元） 27,985 31,629 36,091 42,082 50,980 同比增长率（\%） 22.63\% 13.02\% 14.11\% 16.60\% 21.14\% 归母净利润（百万元） 6,337 7,711 9,613 11,572 14,116 同比增长率（\%） 47.28\% 21.69\% 24.66\% 20.38\% 21.98\% 每股收益（元 / 股） 0.95 1.1……风险段摘要：竞争格局恶化风险、销售不及预期风险、行业政策风险等。 盈利预测与估值（人民币） 2024 2025 2026E 2027E 2028E 营业收入（百万元） 27,985 31,629 36,091 42,082 50,980 同比增长率（\%） 22.63\% 13.02\% 14.11\% 16……

\begin{keyinsight}[AStock处理]
处置为\textbf{正式核心模型}。公司未形成H1预告时序，本轮报告用于Q1和全年预测对照，仍需按当前价格重算。下一步门槛：创新药销售增速不低于30\%、EPS不低于1.35元、关键BD/临床里程碑兑现。
\end{keyinsight}

\section{甘李药业（603087）：两份原始研报对照}

\textbf{中邮证券，2026-05-01，《国内海外销售增长共振，创新管线迈入兑现期》}
元数据给出2026E EPS 2.29元、PE 26.3倍；时序为\textbf{无预告时序}。原文盈利段摘要：预计公司 26 / 27 / 28 年收入分别为 45.62 / 53.89 / 64.91 亿元， 同比增长 12.57\% / 18.13\% / 20.45\%；归母净利润为 13.67 / 16.99 / 21.33 亿元，同比增长 19.55\% / 24.27\% / 25.53\%，对应 PE 为 26 / 21 / 17 倍。 公司是代谢药物龙头，减重药物数据优秀，海外市场增长确定性明显。 首次覆盖，给予“买入”评级。风险段摘要：下游研发需求变动风险；市场竞争加剧风险；地缘政治风险 n 盈利预测和财务指标 项目 / 年度 2025A 2026E 2027E 2028E 营业收入（百万元） 4052 4562 5389 6491 增长率(\%) 33.06 12.57 18.13 20.45 EBITDA（百万元） 1……

\textbf{长城国瑞证券，2026-04-30，《2025年业绩高增，全球化与创新双轮驱动》}
元数据给出2026E EPS 2.28元、PE 26.2倍；时序为\textbf{无预告时序}。原文盈利段摘要：（单位：百万元） 资产负债表 2025 2026E 2027E 2028E 利润表 2025 2026E 2027E 2028E 货币资金 2,088.31 4,011.72 4,656.82 6,181.66 营业收入 4,052.15 4,642.53 5,443.00 6,231.03 应收账款 593.23 309.48 688.40 453.96 营业成本 978.95 1,092.52 1,266.50 1,438.16 预付账款 55.48……风险段摘要：医保降价 / 药品集中招标采购风险、新药研发及商业化过程漫长且成本 高昂风险、研发人员流失风险 主要财务数据及预测： 2025 2026E 2027E 2028E 营业收入(百万元) 4,052.15 4,642.53 5,443.00 6,231.03 增长率(\%) 33.06 14.5……

\begin{keyinsight}[AStock处理]
处置为\textbf{管线观察、等待中报}。公司未形成H1预告时序，本轮报告用于Q1和全年预测对照，仍需按当前价格重算。下一步门槛：等待产品销售、临床/获批、BD里程碑和费用率更新。
\end{keyinsight}

\section{百济神州-U（688235）：两份原始研报对照}

\textbf{西南证券，2026-06-08，《ASCO公布新数据，彰显实体瘤管线加速推进》}
元数据给出2026E EPS 2.85元、PE 78.7倍；时序为\textbf{无预告时序}。原文盈利段摘要：预计 2026-2028 年 EPS 分别为 2.85 / 4.71 / 5.84 元，对应 PE 为 1. 百济神州 （688235）： 生物科技出海明星， 79 / 48 / 38x。 2025 年迎来蝶变 (2025-06-24) 风险提示：临床数据不及预期风险、竞争格局加剧风险、海外拓展不及预期风 险。 指标 / 年度 2025A 2026E 2027E 2028E 营业收入（百万元） 38225.00 45654.23 53536.51 6005……风险段摘要：临床数据不及预期风险、竞争格局加剧风险、海外拓展不及预期风 险。 指标 / 年度 2025A 2026E 2027E 2028E 营业收入（百万元） 38225.00 45654.23 53536.51 60051.52 增长率 40.46\% 19.44\% 17.27\% 12.17\% 归属……

\textbf{中邮证券，2026-05-13，《泽布替尼同比维持高增，全年业绩指引上调》}
元数据给出2026E EPS 2.42元、PE 101.8倍；时序为\textbf{无预告时序}。原文盈利段摘要：预计公司 26 / 27 / 28 年收入分别为 447.5 / 512.3 / 584.2 亿元，同 比增长 17.1\% / 14.5\% / 14.0\%；归母净利润为 37.4 / 54.9 / 71.3 亿元，同 比增长 155.9\% / 46.9\% / 29.8\%，对应 PE 为 102 / 69 / 53 倍。公司为创新 药出海龙头企业，全球商业化能力突出、在研管线即将进入收获期， 维持“买入”评级。风险段摘要：产品销售不及预期；新药临床进度不及预期；新药临床数据不及 预期。 盈利预测和财务指标 项目 / 年度 2025A 2026E 2027E 2028E 营业收入（百万元） 38225 44749 51231 58423 增长率(\%) 40.46 17.07 14.49 14.04 EBITD……

\begin{keyinsight}[AStock处理]
处置为\textbf{管线观察、等待中报}。公司未形成H1预告时序，本轮报告用于Q1和全年预测对照，仍需按当前价格重算。下一步门槛：等待产品销售、临床/获批、BD里程碑和费用率更新。
\end{keyinsight}

\section{中航光电（002179）：两份原始研报对照}

\textbf{东莞证券，2026-03-31，《深度报告：高端互连技术平台，多赛道布局打开成长空间》}
元数据给出2026E EPS 1.21元、PE 28.0倍；时序为\textbf{无预告时序}。原文盈利段摘要：简表（截至 2026 / 3 / 30） 16。风险段摘要：。原材料价格上涨风险、市场需求与订单交付风险、行业竞争 资料来源：ifind，东莞证券研究所 加剧、技术更新迭代风险、合规与监管风险。 本报告的风险等级为中风险。 本报告的信息均来自已公开信息，关于信息的准确性与完整性，建议投资者谨慎判断，据此入市，风险自担。

\textbf{太平洋，2025-11-09，《三季报点评：研发投入持续增长，新业务布局加速推进》}
元数据给出2026E EPS 1.76元、PE 20.1倍；时序为\textbf{无预告时序}。原文盈利段摘要：与投资评级：预计公司 2025-2027 年的净利润为 32.42 亿 元、37.25 亿元、42.07 亿元，EPS 为 1.53 元、1.76 元、1.99 元，对应 PE 为 23 倍、20 倍、18 倍，维持“买入”评级。 风险提示：订单增长不及预期；新业务拓展不及预期。 盈利预测和财务指标 2024 2025E 2026E 2027E 营业收入（百万元） 20685.53 23374.65 26880.84 30375.35 营业收入增长率(\%)……风险段摘要：订单增长不及预期；新业务拓展不及预期。 盈利预测和财务指标 2024 2025E 2026E 2027E 营业收入（百万元） 20685.53 23374.65 26880.84 30375.35 营业收入增长率(\%) 3.04\% 13.00\% 15.00\% 13.00\% 证券 -3.34……

\begin{keyinsight}[AStock处理]
处置为\textbf{订单观察、等待中报}。公司未形成H1预告时序，本轮报告用于Q1和全年预测对照，仍需按当前价格重算。下一步门槛：等待订单、收入确认、合同负债和经营现金流同步改善。
\end{keyinsight}

\section{星网宇达（002829）：两份原始研报对照}

\textbf{东吴证券，2024-04-30，《2023年年报点评：下游需求充裕，随募投项目推进产能释放有望加速》}
元数据给出2026E EPS 1.30元、PE 16.1倍；时序为\textbf{预告前，分母可能失效}。原文盈利段摘要：2022A 2023A 2024E 2025E 2026E 证券 31.83 26.83 22.28 归母净利润（百万元） 215.80 53.98 167.23 217.85 270.03 研究助理 高正泰 执业证书： 同比 33.98 (74.99) 209.82 30.27 23.95 EPS-最新摊薄（元 / 股） 1.04 0.26 0.80 1.05 1.30 P / E（现价\&最新摊薄） 20.09 80.31 25.92 19.90 16.……风险段摘要：1）市场竞争风险；2）经营业绩存在季节性波动；3）下游需 求不及预期。 1 / 3 东吴证券研究所。

\textbf{东吴证券，2023-10-29，《2023年三季报点评：业绩稳步增长，无人系统市场广阔》}
元数据给出2026E EPS ---元、PE ---倍；时序为\textbf{预告前，分母可能失效}。原文盈利段摘要：2022A 2023E 2024E 2025E 执业证书： 营业总收入（百万元） 1,074 1,371 1,698 2,064 同比 40\% 28\% 24\% 22\% 研究助理 高正泰 归属母公司净利润（百万元） 215 307 405 506 执业证书： 同比 34\% 43\% 32\% 25\% 每股收益-最新股本摊薄（元 / 股） 1.04 1.48 1.95 2.44 股价走势 P / E（现价\&最新股本摊薄） 26.54 18.62 14.11 11.……风险段摘要：1）下游需求不及预期；2）募投项目建设不及预期；3）研发 不及预期。 1 / 3 东吴证券研究所。

\begin{keyinsight}[AStock处理]
处置为\textbf{非经常主导、排除}。两份报告均早于最新预告，只能作为旧市场预期；旧EPS不得继续直接入模。下一步门槛：扣非转正前不升级；投资收益不能替代订单、收入和经营现金流。
\end{keyinsight}

\section{中船防务（600685）：两份原始研报对照}

\textbf{华源证券，2025-08-18，《专注综合海洋装备制造，新趋势有望带动业绩增长》}
元数据给出2026E EPS 1.07元、PE 27.1倍；时序为\textbf{预告前，分母可能失效}。原文盈利段摘要：（人民币） 2023 2024 2025E 2026E 2027E 营业收入（百万元） 16,146 19,402 20,489 22,554 23,427 同比增长率（\%） 26.19\% 20.17\% 5.60\% 10.08\% 3.87\% 归母净利润（百万元） 48 377 982 1,513 2,335 同比增长率（\%） -93.02\% 684.86\% 160.20\% 54.18\% 54.25\% 每股收益（元 / 股） 0.03 0.27 0.69 1……风险段摘要：。航运行业景气不及预期；IMO 环保约束不及预期；核心技术风险；地缘 资产负债率（\%） 62.64 政治风险 每股净资产（元 / 股） 12.11 资料来源：聚源数据 盈利预测与估值（人民币） 2023 2024 2025E 2026E 2027E 营业收入（百万元） 16,146 19,……

\textbf{西南证券，2024-04-29，《2024年一季报点评：造船市场火热，利润开始释放》}
元数据给出2026E EPS 1.35元、PE 18.8倍；时序为\textbf{预告前，分母可能失效}。原文盈利段摘要：。预计公司 2024-2026年营业收入分别为 185.7亿元、216.5 1. 中船防务（600685）：新造船市场稳定 亿元和 229.9 亿元，归母净利润分别为 10.2 亿元、14.3 亿元、19 亿元，EPS 向好，公司盈利能力有望提升 分别为 0.72 元、1.01 元、1.35 元。我们认为公司加快生产节奏，持续交付高价 (2023-11-10) 在手订单，毛利有望进一步改善，同时未来看好造船产能有限的大背景下，新 一轮船舶更新周期的开启，将……风险段摘要：汇率波动风险、造船需求不及预期风险、原材料价格人工成本波动 风险。 指标 / 年度 2023A 2024E 2025E 2026E 营业收入（百万元） 16145.95 18568.74 21645.31 22986.22 增长率 26.19\% 15.01\% 16.57\% 6.19\% 归属……

\begin{keyinsight}[AStock处理]
处置为\textbf{预告后重估池}。两份报告均早于最新预告，只能作为旧市场预期；旧EPS不得继续直接入模。下一步门槛：拆分造船主营、联营收益和分红收益，获取预告后2026E分母。
\end{keyinsight}

\section{中航沈飞（600760）：两份原始研报对照}

\textbf{山西证券，2026-04-09，《加速产品转代更迭，内装外贸双轮驱动》}
元数据给出2026E EPS 1.30元、PE 37.5倍；时序为\textbf{无预告时序}。原文盈利段摘要：净资产收益率(\%)： 13.80 资料来源：常闻 预计公司 2026-2028 年 EPS 分别为 1.30 / 1.47 / 1.62，对应公司 4 月 8 日收盘价 48.98 元，2026-2028 年 PE 分别为 37.5 / 33.3 / 30.2，维持“买入-A” 2027E 2028E 李通 营业收入(百万元) 42,837 44,656 47,909 55,305 61,095 执业登记编码： YoY(\%) -7.4 4.2 7.3 15……风险段摘要：执业登记编码： 主力机型列装不及预期；生产交付不及预期。 邮箱： 财务数据与估值： 会计年度 2024A 2025A 2026E 2027E 2028E 李通 营业收入(百万元) 42,837 44,656 47,909 55,305 61,095 执业登记编码： YoY(\%) -7.4……

\textbf{华安证券，2025-11-06，《第三季度经营稳健，军贸或成未来关注方向》}
元数据给出2026E EPS 1.42元、PE 42.6倍；时序为\textbf{无预告时序}。原文盈利段摘要：资产负债表 单位:百万元 利润表 单位:百万元 会计年度 2024 2025E 2026E 2027E 会计年度 2024A 2025E 2026E 2027E 流动资产 A 5002 48129 54682 62661 营业收入 42837 44135 49821 57250 现金 1 9892 16335 18823 21531 营业成本 37497 38552 43515 49943 应收账款 1922 5594 6479 7315 营业税金及附加……风险段摘要：下游需求及交付节奏低于预期；研发不及预期；产能释放不及预期；存 货及应收账款计提风险；公司核心人员变动风险；公司或公司管理层因 触发相关法律法规被罚风险；公司管理层或核心技术人员因无法履行工 作职责影响经营风险。 重要财务指标 单位:百万元 主要财务指标 2024A 2025E 2026E……

\begin{keyinsight}[AStock处理]
处置为\textbf{订单观察、等待中报}。公司未形成H1预告时序，本轮报告用于Q1和全年预测对照，仍需按当前价格重算。下一步门槛：等待订单、收入确认、合同负债和经营现金流同步改善。
\end{keyinsight}

\section{航发动力（600893）：两份原始研报对照}

\textbf{中航证券，2026-06-15，《2025年年报及2026年一季报点评：深化军机民机、国内国际、航机燃机协同布局，奋进第二曲线》}
元数据给出2026E EPS 0.20元、PE 248.3倍；时序为\textbf{无预告时序}。原文盈利段摘要：1联 联联联联联联联联联联联联联联联联联联联联联联联联联联联联联 联联联联联联联联联联联联联联联联联联联联联联联联联联联联联 联联联联联联联联联联联联联联联联联联联联联 联联联联联联联 联联联联联联联联联联联联 2联 联联联联联联联联联联联联联联联联联联联联联联联联联联联联 联联联联联联联联联联联联联联联联联联联联联联联联联联联联联 联联联联联联联联联联联联联联联联联联联联联联联联联联联联联 联联联联联联联联联联联联联联联联联联联联联联联联联联联联联 联联……风险段摘要：联联联联联联联联联联联联联联联联联联联联联联联联联联联联联联 联联联联联联联联联联联联联联联联联联联联联联联联联联联联联联联 联联 联联联联联联联联联联联联联联联联联联联联联联 联联联联联联联联联联联联联联联联联联600893 联2025 联联联联 2026 联联联联联联联联联联联联联联联……

\textbf{国信证券，2026-04-14，《“航改燃+大飞机+军品后装”三轮驱动蓄势待发》}
元数据给出2026E EPS 0.20元、PE 262.1倍；时序为\textbf{无预告时序}。原文盈利段摘要：和财务指标 2024 2025 2026E 2027E 2028E 营业收入(百万元) 47,880 46,331 50,038 55,042 60,546 (+ / -\%) 9.5\% -3.2\% 8.0\% 10.0\% 10.0\% 归母净利润(百万元) 860 634 523 782 967 (+ / -\%) -39.5\% -26.3\% -17.5\% 49.3\% 23.7\% 每股收益（元） 0.32 0.24 0.20 0.29 0.36 EBIT Mar……风险段摘要：技术迭代不及预期；市场需求不及预期；传统主业竞争加剧。 投资建议：首次覆盖，给予“优于大市”评级。综合绝对估值与相对估值， 我们认为公司股票合理价值在 55.56-59.28 元之间，相对于公司目前股 价有 8.0\%-15.2\%溢价空间。我们预计公司 2026-2028 年归母净利润为 5……

\begin{keyinsight}[AStock处理]
处置为\textbf{订单观察、等待中报}。公司未形成H1预告时序，本轮报告用于Q1和全年预测对照，仍需按当前价格重算。下一步门槛：等待订单、收入确认、合同负债和经营现金流同步改善。
\end{keyinsight}

\section{徐工机械（000425）：两份原始研报对照}

\textbf{华鑫证券，2026-05-27，《公司动态研究报告：营收稳健增长，国际化与矿机双轮驱动开启新周期》}
元数据给出2026E EPS 0.66元、PE 14.5倍；时序为\textbf{无预告时序}。原文盈利段摘要：预测公司 2026-2028 年收入分别为 1147.77、1294.68、1443.57 亿 元，EPS 分别为 0.66、0.83、0.97 元， 周期波动风险；海外市场地缘政治及汇率波动风险；原 材料价格波动导致毛利率修复不及预期风险；矿业机械放量节奏不及 预期风险；行业竞争加剧导致价格体系承压风险。。 预测指标 2025A 2026E 2027E 2028E 主营收入（百万元） 100,823 114,777 129,468 144,357 增长率……风险段摘要：工程机械行业周期波动风险；海外市场地缘政治及汇率波动风险；原 材料价格波动导致毛利率修复不及预期风险；矿业机械放量节奏不及 预期风险；行业竞争加剧导致价格体系承压风险。。 预测指标 2025A 2026E 2027E 2028E 主营收入（百万元） 100,823 114,777 129,……

\textbf{中邮证券，2026-05-26，《业绩稳健增长，现金流大幅改善》}
元数据给出2026E EPS 0.67元、PE 14.4倍；时序为\textbf{无预告时序}。原文盈利段摘要：预计公司 2026-2028 年营收分别为 1141.63、1276.14、1408.54 亿元，同比增速分别为 13.23\%、11.78\%、10.37\%；归母净利润分别为 78.31、99.71、116.58 亿元，同比增速分别为 19.16\%、27.33\%、16.92\%。 公司 2026-2028 年业绩对应 PE 估值分别为 14.43、11.33、9.69，维 持“买入”评级。 l 风险提示： 行业景气度不及预期风险；海外市场开拓不及预期风险；竞争……风险段摘要：行业景气度不及预期风险；海外市场开拓不及预期风险；竞争加。

\begin{keyinsight}[AStock处理]
处置为\textbf{正式核心模型}。公司未形成H1预告时序，本轮报告用于Q1和全年预测对照，仍需按当前价格重算。下一步门槛：海外和矿机收入双位数增长、经营现金流持续为正、EPS不低于0.60元。
\end{keyinsight}

\section{渝农商行（601077）：两份原始研报对照}

\textbf{开源证券，2026-06-08，《公司首次覆盖报告：高股息筑底、量增价稳释放业绩弹性》}
元数据给出2026E EPS 1.13元、PE 5.9倍；时序为\textbf{无预告时序}。原文盈利段摘要：原文目录未形成可稳定抽取的盈利段，保留EPS元数据与PDF路径供复核。风险段摘要：经济增速不及预期；政策落地不及预期；业务转型的风险等。 财务摘要和估值指标 指标 2024A 2025A 2026E 2027E 2028E 营业收入(百万元) 28,261 28,648 29,289 30,707 31,952 YOY(\%) 1.09 1.37 2.24 4.84 4.……

\textbf{太平洋，2026-04-29，《渝农商行2025年年报点评：扩表快、息差稳》}
元数据给出2026E EPS ---元、PE ---倍；时序为\textbf{无预告时序}。原文盈利段摘要：公司全力推进“三种新动能”战略，强化核心竞争力；在 、343.95 亿元，归母净利润为 131.06、142.75、156.02 亿元，每股净资产为 12.16、 13.30、14.19 元，对应 4 月 27 日收盘价的 PB 估值为 0.61、0.56、0.52 倍。维持“买入”评级。 风险提示：经济复苏不及预期、息差持续收窄、资产质量恶化。风险段摘要：经济复苏不及预期、息差持续收窄、资产质量恶化。

\begin{keyinsight}[AStock处理]
处置为\textbf{正式核心模型}。公司未形成H1预告时序，本轮报告用于Q1和全年预测对照，仍需按当前价格重算。下一步门槛：净息差不低于1.60\%、不良率不高于1.10\%，营收维持正增长。
\end{keyinsight}

\section{沪农商行（601825）：两份原始研报对照}

\textbf{国信证券，2026-04-26，《2025年报及2026一季报点评：息差降幅收窄，资产质量稳中向好》}
元数据给出2026E EPS 1.29元、PE 7.1倍；时序为\textbf{无预告时序}。原文盈利段摘要：和财务指标 2024 2025 2026E 2027E 2028E 营业收入(百万元) 26,641 25,870 26,000 27,215 28,421 (+ / -\%) 0.9\% -2.9\% 0.5\% 4.7\% 4.4\% 归母净利润(百万元) 12,288 12,313 12,399 12,625 12,822 (+ / -\%) 1.2\% 0.2\% 0.7\% 1.8\% 1.6\% 摊薄每股收益(元) 1.27 1.28 1.29 1.31 1.33 总……风险段摘要：宏观经济形势走弱可能对银行资产质量产生不利影响。 盈利预测和财务指标 2024 2025 2026E 2027E 2028E 营业收入(百万元) 26,641 25,870 26,000 27,215 28,421 (+ / -\%) 0.9\% -2.9\% 0.5\% 4.7\% 4.4\% 归母……

\textbf{国信证券，2025-09-01，《2025年中报点评：营收降幅收窄，利润小幅增长》}
元数据给出2026E EPS 1.32元、PE 6.6倍；时序为\textbf{无预告时序}。原文盈利段摘要：预计公司 2025-2027 年归 母净利润 124 / 127 / 132 亿元（上次预测值 124 / 128 / 132 亿元），同比增速 0.6 / 2.7 / 4.1\% ； 摊 薄 EPS 为 1.28 / 1.32 / 1.37 元 ； 当 前 股 价 对 应 PE 为 6.8 / 6.6 / 6.4x，PB 为 0.64 / 0.60 / 0.56x，维持“中性”评级。 风险提示：宏观经济形势走弱可能对银行资产质量产生不利影响。 盈利预测和财……风险段摘要：宏观经济形势走弱可能对银行资产质量产生不利影响。 盈利预测和财务指标 2023 2024 2025E 2026E 2027E 营业收入(百万元) 26,414 26,641 25,935 26,410 27,691 (+ / -\%) 3.1\% 0.9\% -2.7\% 1.8\% 4.8\% 归母……

\begin{keyinsight}[AStock处理]
处置为\textbf{正式核心模型}。公司未形成H1预告时序，本轮报告用于Q1和全年预测对照，仍需按当前价格重算。下一步门槛：营收维持正增长、拨备覆盖率不低于300\%，分红率稳定。
\end{keyinsight}
\begin{riskbox}[研报使用边界]
券商EPS、PE和目标价是市场预期锚，不是AStock的事实判断。预告前报告可以证明市场此前预期有多低，但不能在预告后继续充当当前分母。仅有“买入/强推”而无公开目标价时，Street目标权重必须为0。
\end{riskbox}
